\documentclass[12pt]{article}
\usepackage[margin=1in]{geometry}
\usepackage{hyperref}

\hypersetup{colorlinks=true,linkcolor=blue,urlcolor=blue}

\begin{document}

\begin{flushright}
Yaoping Wang\\
Department of Statistics, Rutgers University\\
Piscataway, NJ, USA\\
yw1580@scarletmail.rutgers.edu\\
\end{flushright}

\vspace{2em}

\begin{flushleft}
Editor-in-Chief\\
\textit{International Journal of Data Science and Analytics}\\
Springer
\end{flushleft}

\vspace{1em}

Dear Editor,

\vspace{1em}

I am submitting the manuscript \textbf{"Protocol-Dependent Resolution in Finite-Sample Model Selection"} for consideration as a regular paper in the \textit{International Journal of Data Science and Analytics}.

\paragraph{The finding.}
A single decoding parameter, the maximum number of tokens a model can generate, can reverse whether a set of candidate models looks distinguishable or not. It can also hide a negative trend behind what appears to be flat performance. None of 50 recent fine-tuning papers report this parameter. The paper documents the finding through LoRA experiments on two model families across several training seeds.

\paragraph{The diagnostic.}
The threshold I use is $\Delta_{\min}$, the minimum detectable effect at 50\% power for a two-proportion $z$-test. The formula itself is textbook power analysis. What matters is what applying it reveals: evaluation resolution depends on the protocol, a short generation length compresses between-checkpoint variance and creates false negatives, and model comparison claims are not interpretable without a fully specified evaluation pipeline.

\paragraph{The evidence.}
Four experiments support this. First, moving from 128 to 256 tokens shifts observed accuracy ranges by $2.3\times$ and flips the pairwise resolution verdict across three seeds. Between-checkpoint standard deviation more than doubles. Second, the longer token limit reveals a statistically significant negative trend in one seed ($\rho=-0.77$, $p=0.009$) that is invisible under truncation. Third, Mistral-7B, which generates shorter outputs, shows much less sensitivity to the token limit, which confirms that truncated chain-of-thought reasoning is the mechanism. Fourth, all three seeds are close to the Bonferroni-corrected all-pairwise threshold. At $N=200$, even a fully specified protocol cannot support selecting the best from ten candidates. The evaluation budget itself is the binding constraint.

\paragraph{Honesty about limits.}
I report a case where the diagnostic fails silently (near-floor accuracy), document nine pitfalls I ran into during the research, and note that the soup gain ($\kappa \approx +2.0$~pp) is measured on one trajectory. A survey of 50 recent fine-tuning papers found that none report the token limit parameter that our experiments show can determine the verdict. A diagnostic that never reports when it does not apply is not very useful, so I delimit the framework rather than oversell it.

\paragraph{Fit with JDSA.}
The work fits within the journal's scope of statistical foundations for data science. The diagnostic is model- and metric-agnostic, and the core finding, that an unreported protocol parameter can reverse empirical conclusions, generalises beyond any one model class. A practitioner quick-reference table and an open-source implementation are included.

\paragraph{Declarations.}
The manuscript is original, has not been published before, and is not under consideration elsewhere. I am the sole author. I have no competing interests and received no funding for this work. All code and data are openly available at \url{https://github.com/wyp629929/A-Resolution-Framework-for-Finite-Sample-Model-Selection}.

\vspace{2em}

Thank you for your time and consideration.

\vspace{1em}

Sincerely,

\vspace{2em}

Yaoping Wang

\end{document}
